\subsection[Proof of Theorem~\ref{thm:multibaileylattice}]{Proof of Theorem~\ref{thm:multibaileylattice}}

Now we can turn to the proof of Theorem~\ref{thm:multibaileylattice}.
Recall the classical $q$-analogues of Pascal's triangle
\begin{equation}\label{eq:pascal1}
\left[{N+1\atop j}\right] = q^j\left[{N\atop j}\right]+\left[{N\atop j-1}\right],
\end{equation}
and
\begin{equation} \label{eq:pascal2}
\left[{N+1\atop j}\right]=\left[{N\atop j}\right]+ q^{N+1-j}\left[{N\atop j-1}\right],
\end{equation}
for all integers $N$, $j$ with $N\geq 0$.

\subsubsection[Recurrence relation for f\_\{N,j,n\}(b\_1,...,b\_N)]{Recurrence relation for $\boldsymbol{f_{N,j,n}(b_1,\ldots,b_N)}$}
We first prove the following recurrence relation, which plays a central role in our proof.
\begin{Proposition}[Recurrence relation]\label{prop:rec} For all $0\leq j\leq N+1$,
\begin{gather*}
\bigl(1-aq^{2n-1-N}\bigr)f_{N+1,j,n}(b_1,\ldots,b_{N+1})\\
\qquad =(1-b_{N+1}q^n)\bigl(1-aq^{2n-1-N-j}\bigr)f_{N,j,n}(b_1,\ldots,b_N)\\
\phantom{\qquad =}{}+\bigl(aq^{n-1-N}-b_{N+1}\bigr)\bigl(1-aq^{2n-j}\bigr)f_{N,j-1,n-1}(b_1,\ldots,b_N).
\end{gather*}
\end{Proposition}

We start by giving two technical lemmas which are extensions of \eqref{eq:pascal1} and \eqref{eq:pascal2}. Their proofs are straightforward verifications and are therefore omitted here.


\begin{Lemma}
\label{lem:tech1}
For all $j,u \in \Z$ and $0\leq M\leq N$,
\begin{gather*}
\begin{bmatrix}
M\\
j-u
\end{bmatrix} \begin{bmatrix}
N+1-M\\
u
\end{bmatrix}-q^{2j-2u-M+1}\begin{bmatrix}
M\\
j-u+1
\end{bmatrix} \begin{bmatrix}
N+1-M\\
u-1
\end{bmatrix}\\
\qquad=q^{u}\begin{bmatrix}
M\\
j-u
\end{bmatrix} \begin{bmatrix}
N-M\\
u
\end{bmatrix}
+q^{j-u-M}\begin{bmatrix}
N-M\\
u-1
\end{bmatrix}\left(\begin{bmatrix}
M\\
j-u
\end{bmatrix}-\begin{bmatrix}
M\\
j-u+1
\end{bmatrix}\right)\\
\phantom{\qquad=}{} -q^{2j-3u-2M+3+N}\begin{bmatrix}
M\\
j-u+1
\end{bmatrix} \begin{bmatrix}
N-M\\
u-2
\end{bmatrix} .
\end{gather*}
\end{Lemma}

Note that when $M=0$ and $u=j$, Lemma~\ref{lem:tech1} reduces to \eqref{eq:pascal1}.

\begin{Lemma}
\label{lem:tech2}
For all $j,u \in \Z$ and $0\leq M\leq N$,
\begin{gather*}
\begin{bmatrix}
M+1\\
j-u
\end{bmatrix} \begin{bmatrix}
N-M\\
u
\end{bmatrix}-q^{2j-2u-M}\begin{bmatrix}
M+1\\
j-u+1
\end{bmatrix} \begin{bmatrix}
N-M\\
u-1
\end{bmatrix}\\
\qquad
=q^j \begin{bmatrix}
M\\
j-u
\end{bmatrix} \begin{bmatrix}
N-M\\
u
\end{bmatrix}
- q^{2j-2u-M}\begin{bmatrix}
M\\
j-u+1
\end{bmatrix} \begin{bmatrix}
N-M\\
u-1
\end{bmatrix}+\begin{bmatrix}
M\\
j-u-1
\end{bmatrix} \begin{bmatrix}
N-M\\
u
\end{bmatrix}\\
\phantom{\qquad=}{}- q^{N+j+1-2u-M} \begin{bmatrix}
M\\
j-u
\end{bmatrix} \begin{bmatrix}
N-M\\
u-1
\end{bmatrix}.
\end{gather*}
\end{Lemma}

Note that when $M=1$ and $u=j$, Lemma~\ref{lem:tech2} reduces to a combination of \eqref{eq:pascal1} and \eqref{eq:pascal2}.

We can now prove our recurrence relation.
\begin{proof}[Proof of Proposition \ref{prop:rec}]
Using the homogeneity of $e_M$, recall from Theorem~\ref{thm:multibaileylattice} that
\begin{gather*}
f_{N,j,n}(b_1,\ldots,b_N)\\
\qquad= \sum_{M \in \Z} \sum_{u \in \Z} a^{u}q^{(M-j+u)(n-j+u)+u(n-N)} \begin{bmatrix}
M\\
j-u
\end{bmatrix} \begin{bmatrix}
N-M\\
u
\end{bmatrix} e_M(-b_1, \dots , -b_N).
\end{gather*}
Thus
\begin{gather}
\bigl(1-aq^{2n-1-N}\bigr)f_{N+1,j,n}(b_1,\ldots,b_{N+1})\nonumber \\
\qquad= \sum_{M \in \Z} \sum_{u\in \mathbb{Z}} a^{u}q^{(M-j+u)(n-j+u)+u(n-N-1)}\nonumber\\
\phantom{\qquad=}{}\times\bigg( \begin{bmatrix}
M\\
j-u
\end{bmatrix} \begin{bmatrix}
N+1-M\\
u
\end{bmatrix}
-q^{2j-2u-M+1}\begin{bmatrix}
M\\
j-u+1
\end{bmatrix} \begin{bmatrix}
N+1-M\\
u-1
\end{bmatrix} \bigg)\nonumber\\
\phantom{\qquad=}{}\times e_M(-b_1, \dots , -b_{N+1}).\label{eq:lastline}
\end{gather}
Writing
\begin{equation*}e_M(-b_1, \dots , -b_{N+1})=e_M(-b_1, \dots , -b_{N})-b_{N+1}e_{M-1}(-b_1, \dots , -b_{N}),
\end{equation*}
it is enough to show that coefficients of $b_{N+1}^0a^ue_M(-b_1, \dots , -b_{N})$ and $b_{N+1}^1a^ue_M(-b_1, \dots , -b_{N})$ coincide on both sides of Proposition~\ref{prop:rec}.

First, we derive from~\eqref{eq:lastline} that the coefficient of $b_{N+1}^0a^ue_M(-b_1, \dots , -b_{N})$ on the left-hand side of Proposition~\ref{prop:rec} is equal to \smash{$q^{(M-j+u)(n-j+u)+u(n-N-1)}$} times
\begin{equation}\label{eq:coeffb0l}
\begin{bmatrix}
M\\
j-u
\end{bmatrix} \begin{bmatrix}
N+1-M\\
u
\end{bmatrix}-q^{2j-2u-M+1}\begin{bmatrix}
M\\
j-u+1
\end{bmatrix} \begin{bmatrix}
N+1-M\\
u-1
\end{bmatrix}.
\end{equation}

Now the coefficient of $b_{N+1}^0a^ue_M(-b_1, \dots , -b_{N})$ on the right-hand side of Proposition~\ref{prop:rec} is equal to \smash{$q^{(M-j+u)(n-j+u)+u(n-N-1)}$} times
\begin{gather}
q^u \begin{bmatrix}
M\\
j-u
\end{bmatrix} \begin{bmatrix}
N-M\\
u
\end{bmatrix}-q^{j-u-M}\begin{bmatrix}
M\\
j-u+1
\end{bmatrix} \begin{bmatrix}
N-M\\
u-1
\end{bmatrix}+q^{j-u-M}\begin{bmatrix}
M\\
j-u
\end{bmatrix} \begin{bmatrix}
N-M\\
u-1
\end{bmatrix}\nonumber\\
\qquad{}-q^{N+2j-2M-3u+3} \begin{bmatrix}
M\\
j-u+1
\end{bmatrix} \begin{bmatrix}
N-M\\
u-2
\end{bmatrix}.\label{eq:coeffb0r}
\end{gather}
Note finally that~\eqref{eq:coeffb0l} (resp.~\eqref{eq:coeffb0r}) is the left-hand (resp.\ right-hand) side of Lemma~\ref{lem:tech1}, so we are done regarding this first coefficient.

Similarly, the coefficient of $b_{N+1}^1a^ue_M(-b_1, \dots , \!-b_{N})$ in~\eqref{eq:lastline} is \smash{$-q^{(M+1-j+u)(n-j+u)+u(n-N-1)}$} times
\begin{equation}\label{eq:coeffb1l}
\begin{bmatrix}
M+1\\
j-u
\end{bmatrix} \begin{bmatrix}
N-M\\
u
\end{bmatrix}-q^{2j-2u-M}\begin{bmatrix}
M+1\\
j-u+1
\end{bmatrix} \begin{bmatrix}
N-M\\
u-1
\end{bmatrix},
\end{equation}

and the coefficient of $b_{N+1}^1a^ue_M(-b_1, \dots , -b_{N})$ on the right-hand side of Proposition~\ref{prop:rec} is equal to \smash{$-q^{(M+1-j+u)(n-j+u)+u(n-N-1)}$} times
\begin{gather}
q^j \begin{bmatrix}
M\\
j-u
\end{bmatrix} \begin{bmatrix}
N-M\\
u
\end{bmatrix}-q^{2j-2u-M}\begin{bmatrix}
M\\
j-u+1
\end{bmatrix} \begin{bmatrix}
N-M\\
u-1
\end{bmatrix}+\begin{bmatrix}
M\\
j-u-1
\end{bmatrix} \begin{bmatrix}
N-M\\
u
\end{bmatrix}\nonumber\\
\qquad{}-q^{N+j+1-2u-M} \begin{bmatrix}
M\\
j-u
\end{bmatrix} \begin{bmatrix}
N-M\\
u-1
\end{bmatrix}.\label{eq:coeffb1r}
\end{gather}

As~\eqref{eq:coeffb1l} (resp.~\eqref{eq:coeffb1r}) is the left-hand (resp.\ right-hand) side of Lemma~\ref{lem:tech2}, this ends the proof of Proposition~\ref{prop:rec}.
\end{proof}

\subsubsection[Proof of Theorem \ref{thm:multibaileylattice}]{Proof of Theorem \ref{thm:multibaileylattice}}

Now we use Proposition \ref{prop:rec} and Lemma \ref{lem:baileylattices} to prove Theorem \ref{thm:multibaileylattice}.

We proceed by induction on $N$. For $N=0$, by \eqref{eq:multibaileylatticebeta}, \smash{$\beta^{(0)}_n=\beta_n$} and by \eqref{eq:multibaileylatticealpha},
\begin{equation*}
\alpha^{(0)}_n = \bigl(1-aq^{2n}\bigr)\frac{\alpha_{n}}{1-aq^{2n}}=\alpha_n.
\end{equation*}
Thus $\bigl(\alpha^{(0)}_n, \beta^{(0)}_n\bigr)$ is a bilateral Bailey pair relative to $a$.

Now, assume that for some integer $N\geq 0$, \smash{$\bigl(\alpha^{(N)}_n,\beta^{(N)}_n\bigr)$} is a bilateral Bailey pair relative to~$aq^{-N}$ and show that \smash{$\bigl(\alpha^{(N+1)}_n,\beta^{(N+1)}_n\bigr)$} is a bilateral Bailey pair relative to $aq^{-N-1}$. By Lemma~\ref{lem:baileylattices} with $b = b_{N+1}$, $(\alpha'_n,\beta'_n)$ is a bilateral Bailey pair relative to $aq^{-N-1}$, where
\begin{equation*}
\alpha'_n=\frac{1-aq^{-N}}{1-b_{N+1}}\left(\frac{(1-b_{N+1}q^n) \alpha_n^{(N)}}{1-aq^{2n-N}}-\frac{q^{n-1}\bigl(aq^{n-N-1}-b_{N+1}\bigr) \alpha_{n-1}^{(N)}}{1-aq^{2n-N-2}}\right),
\end{equation*}
and
\begin{equation*}\beta'_n=\frac{1-b_{N+1}q^n}{1-b_{N+1}}\beta_n^{(N)}.
\end{equation*}

Now let us show that $(\alpha'_n,\beta'_n) = \bigl(\alpha^{(N+1)}_n,\beta^{(N+1)}_n\bigr)$.
It is clear that $\beta'_n=\beta^{(N+1)}_n$, and for $\alpha'_n$, by~\eqref{eq:multibaileylatticealpha} we have
\begin{align*}
\alpha'_n
={}&\frac{1-aq^{-N}}{1-b_{N+1}}\left(\frac{(1-b_{N+1}q^n) \alpha_n^{(N)}}{1-aq^{2n-N}}-\frac{q^{n-1}\bigl(aq^{n-N-1}-b_{N+1}\bigr) \alpha_{n-1}^{(N)}}{1-aq^{2n-N-2}}\right)\\
={}& \frac{\bigl(aq^{-N}\bigr)_{N+1}}{(1-b_1)\cdots(1-b_{N+1})} \bigg(\sum_{j \in \Z}(-1)^j\frac{q^{jn-j(j+1)/2}(1-b_{N+1}q^n)f_{N,j,n}(b_1,\ldots,b_N)}{\bigl(aq^{2n-N-j}\bigr)_{N+1}} \alpha_{n-j} \\
& +\sum_{j \in \Z}(-1)^{j+1}\frac{q^{(j+1)(n-1)-j(j+1)/2}\bigl(aq^{n-1-N}-b_{N+1}\bigr)f_{N,j,n-1}(b_1,\ldots,b_N)}{\bigl(aq^{2n-2-N-j}\bigr)_{N+1}} \alpha_{n-1-j} \bigg)\\
={}& \frac{\bigl(aq^{-N}\bigr)_{N+1}}{(1-b_1)\cdots(1-b_{N+1})} \sum_{j \in \Z}(-1)^j\frac{q^{jn-j(j+1)/2}}{(aq^{2n-N-j-1})_{N+2}} \big( (1-b_{N+1}q^n)\bigl(1-aq^{2n-N-j-1}\bigr)
\\
& \times f_{N,j,n}(b_1,\ldots,b_N)+ \bigl(aq^{n-1-N}-b_{N+1}\bigr)\bigl(1-aq^{2n-j}\bigr)f_{N,j-1,n-1}(b_1,\ldots,b_N) \big) \alpha_{n-j}.
\end{align*}
Now by Proposition \ref{prop:rec}, this equals
\begin{align*}
\alpha'_n
={}& \frac{\bigl(aq^{-N}\bigr)_{N+1}}{(1-b_1)\cdots(1-b_{N+1})} \sum_{j \in \Z}(-1)^j\frac{q^{jn-j(j+1)/2}}{\bigl(aq^{2n-N-j-1}\bigr)_{N+2}} \bigl(1-aq^{2n-1-N}\bigr)\\
&\times f_{N+1,j,n}(b_1,\ldots,b_{N+1}) \alpha_{n-j}
= \alpha_n^{(N+1)}.
\end{align*}
Thus the pair \smash{$\bigl(\alpha^{(N+1)}_n, \beta^{(N+1)}_n\bigr)$} is indeed a Bailey pair relative to $aq^{-N-1}$.
